% ═══════════════════════════════════════════════════════════════════════
% QOU Formalization Blueprint — content.tex
%
% This file maps every theorem-like object in the manuscript to its
% Lean 4 declaration via \lean{} macros and records dependency edges
% via \uses{}.  Objects whose Lean proof is sorry-free carry \leanok.
%
% Convention:
%   - Labels match the manuscript's \label{} tags exactly.
%   - Lean declarations live in the QOU.* namespace.
%   - \uses{} lists the labels this object depends on.
% ═══════════════════════════════════════════════════════════════════════

\chapter{Core Algebraic Structures (Glossary)}

% ── Frobenius-Hopf Object ──────────────────────────────────────────
\begin{definition}[Frobenius-Hopf Object]
  \label{def:frobenius-hopf}
  \lean{QOU.Glossary.FrobeniusHopfObject}
  An object $A$ in a rigid monoidal category equipped with a Hopf
  algebra structure $(\mu, \eta, \Delta, \varepsilon, S)$ and a
  Frobenius form $\lambda : A \to \mathbf{1}$ satisfying non-degeneracy.
\end{definition}

% ── Quantum Deformation Parameter ─────────────────────────────────
\begin{definition}[Quantum Deformation Parameter]
  \label{def:quantum-parameter}
  \lean{QOU.Glossary.quantumParameter}
  \uses{def:frobenius-hopf}
  The parameter $q \in k^\times$ governing the Hopf algebra structure.
  At $q = 1$ we recover the classical cocommutative case.
\end{definition}

% ── Fibre Functor ──────────────────────────────────────────────────
\begin{definition}[Fibre Functor]
  \label{def:fibre-functor}
  \lean{QOU.Glossary.FibreFunctor}
  An exact faithful $k$-linear monoidal functor
  $\omega : \mathcal{C} \to \mathrm{Vect}_k$ used in Tannaka-Krein
  reconstruction.
\end{definition}

% ── Quantum Gauge Group ───────────────────────────────────────────
\begin{definition}[Quantum Gauge Group]
  \label{def:quantum-gauge-group}
  \lean{QOU.Glossary.QuantumGaugeGroup}
  \uses{def:fibre-functor, def:quantum-parameter}
  The quantum group $G_q = \mathrm{Aut}^{\otimes}(\omega)$ reconstructed
  via Tannaka-Krein duality from the fibre functor.
\end{definition}

% ── CR Manifold ────────────────────────────────────────────────────
\begin{definition}[CR Manifold Structure]
  \label{def:cr-manifold}
  \lean{QOU.Glossary.CRManifold}
  A pseudo-Hermitian structure with Levi distribution, Reeb vector
  field, and Webster metric.
\end{definition}

% ── Bring's Surface ───────────────────────────────────────────────
\begin{definition}[Bring's Surface]
  \label{def:brings-surface}
  \lean{QOU.Glossary.BringsSurface}
  \uses{def:cr-manifold}
  The maximally symmetric CR manifold of genus 4, with Euler
  characteristic $\chi = -6$, defined by $\sum x_j^k = 0$ for
  $k = 1,2,3$ in $\mathbb{CP}^4$.
\end{definition}

% ── Substrate Parameter ───────────────────────────────────────────
\begin{definition}[Substrate Parameter]
  \label{def:substrate-parameter}
  \lean{QOU.Glossary.substrateParameter}
  \uses{def:brings-surface, def:quantum-parameter}
  $q(p) = 1 + |\langle L_\theta, F_A^+\rangle|(p)$ --- couples
  chiral curvature to contact geometry on Bring's surface.
\end{definition}

% ── Jet Prolongation ──────────────────────────────────────────────
\begin{definition}[Jet Prolongation]
  \label{def:jet-prolongation}
  \lean{QOU.Glossary.JetProlongation}
  The first jet bundle $J^1 E$ represented categorically by
  $A \otimes \Omega^1(A)$.
\end{definition}

% ── Kähler Differentials ──────────────────────────────────────────
\begin{definition}[Kähler Differentials]
  \label{def:kahler-differentials}
  \lean{QOU.Glossary.KahlerDifferentials}
  \uses{def:jet-prolongation}
  $\Omega^1(A) = \ker(m : A \otimes A \to A)$, the universal
  first-order differential calculus.
\end{definition}

% ── Braid Group ───────────────────────────────────────────────────
\begin{definition}[Braid Group $B_3$]
  \label{def:braid-group}
  \lean{QOU.Glossary.BraidGroup}
  The braid group on 3 strands: $\langle \sigma_1, \sigma_2 \mid
  \sigma_1\sigma_2\sigma_1 = \sigma_2\sigma_1\sigma_2 \rangle$.
\end{definition}

% ── Binary Dihedral Group ─────────────────────────────────────────
\begin{definition}[Binary Dihedral Group]
  \label{def:binary-dihedral}
  \lean{QOU.Glossary.BinaryDihedralGroup}
  \uses{def:braid-group}
  $\mathrm{Dic}_3 \subset \mathrm{SU}(2)$ of order 12, lifting $S_3$
  via the double cover $\mathrm{SU}(2) \to \mathrm{SO}(3)$.
\end{definition}

% ── F-Isocrystal ──────────────────────────────────────────────────
\begin{definition}[F-Isocrystal]
  \label{def:f-isocrystal}
  \lean{QOU.Glossary.FIsocrystal}
  A module over the Dieudonné ring $\mathbb{Q}_p\langle F, V\rangle
  / (FV - p)$.
\end{definition}


\chapter{Knot Theory and Combinatorial Types}

% ── Planar Diagram ────────────────────────────────────────────────
\begin{definition}[Planar Diagram Code]
  \label{def:planar-diagram}
  \lean{QOU.KnotTheory.PlanarDiagram}
  A combinatorial representation of a knot via a list of crossings,
  each carrying four strand indices and a sign.
\end{definition}

% ── Unknot ────────────────────────────────────────────────────────
\begin{definition}[Unknot ($0_1$)]
  \label{def:unknot}
  \lean{QOU.KnotRegistry.unknotEntry}
  \leanok
  \uses{def:planar-diagram}
  The trivial knot with 0 crossings.  QOU identity: Neutrino (SU(2) vacuum).
  \textit{Knot Atlas}: \url{http://katlas.org/wiki/0_1}
\end{definition}

% ── Trefoil ───────────────────────────────────────────────────────
\begin{definition}[Trefoil ($3_1$)]
  \label{def:trefoil}
  \lean{QOU.KnotRegistry.trefoilEntry}
  \leanok
  \uses{def:planar-diagram, def:braid-group}
  The simplest non-trivial knot (3 crossings, torus knot).
  QOU identity: Electron (Generation~1), source of the $B_3$ monodromy
  for Bring's Surface.
  \textit{Knot Atlas}: \url{http://katlas.org/wiki/3_1}
\end{definition}

% ── Figure-Eight ──────────────────────────────────────────────────
\begin{definition}[Figure-Eight Knot ($4_1$)]
  \label{def:figure-eight}
  \lean{QOU.KnotRegistry.figureEightEntry}
  \leanok
  \uses{def:planar-diagram}
  The simplest hyperbolic knot (4 crossings).
  QOU identity: Muon (Generation~2).  Its hyperbolic volume
  $\mathrm{Vol}(4_1) \approx 2.0298832128$ drives the mass derivation.
  \textit{Knot Atlas}: \url{http://katlas.org/wiki/4_1}
\end{definition}

% ── Hyperbolic Volume ─────────────────────────────────────────────
\begin{definition}[Hyperbolic Volume]
  \label{def:hyperbolic-volume}
  \lean{QOU.KnotTheory.HyperbolicVolume}
  The volume of the hyperbolic knot complement $S^3 \setminus K$,
  computed via the Bloch-Wigner dilogarithm (SnapPy census).
\end{definition}

% ── Vol(4_1) ──────────────────────────────────────────────────────
\begin{definition}[Figure-Eight Volume]
  \label{def:vol-figure-eight}
  \lean{QOU.KnotTheory.vol_figure_eight}
  \leanok
  \uses{def:hyperbolic-volume, def:figure-eight}
  $\mathrm{Vol}(4_1) = 2.0298832128 \pm 10^{-10}$.
  Verified against SnapPy and the Bloch-Wigner function.
\end{definition}


\chapter{Topological Mass Derivation}

% ── Topological Mass Conjecture ───────────────────────────────────
\begin{conjecture}[Topological Decomposition of Mass]
  \label{conj:topological-mass}
  \lean{QOU.MassDerivation.topological_mass_conjecture}
  \uses{def:hyperbolic-volume, def:substrate-parameter}
  For a particle identified with knot $K$, its mass satisfies
  $m = \mathrm{Vol}(K) / \hbar_q^2$ where $\hbar_q = 1 - 1/q$.
  This is Conjecture~4.25 of the manuscript, assumed as an axiom.
\end{conjecture}

% ── CODATA Mass Ratio ─────────────────────────────────────────────
\begin{definition}[Experimental Muon-Electron Mass Ratio]
  \label{def:mu-e-ratio}
  \lean{QOU.MassDerivation.experimental_mu_e_ratio}
  \leanok
  $m_\mu / m_e = 206.768283$ (CODATA 2022).
\end{definition}

% ── Quantum Planck Constant ───────────────────────────────────────
\begin{definition}[Quantum Planck Constant]
  \label{def:quantum-planck}
  \lean{QOU.MassDerivation.quantum_planck}
  \leanok
  \uses{def:quantum-parameter}
  $\hbar_q = 1 - 1/q$ for $q > 1$.
\end{definition}

% ── Substrate Coupling Constant ───────────────────────────────────
\begin{definition}[Substrate Coupling Constant]
  \label{def:substrate-q}
  \lean{QOU.MassDerivation.substrate_q}
  \leanok
  $q \approx 1.1097$, the derived value of the substrate parameter.
\end{definition}

% ── Vol Approximation ─────────────────────────────────────────────
\begin{lemma}[Figure-Eight Volume Approximation]
  \label{lem:vol-approx}
  \lean{QOU.MassDerivation.vol_figure_eight_approx}
  \uses{def:vol-figure-eight}
  $|\mathrm{Vol}(4_1) - 2.02988| < 10^{-5}$.
\end{lemma}

% ── Mass Ratio Approximation ──────────────────────────────────────
\begin{lemma}[Mass Ratio Approximation]
  \label{lem:mu-e-approx}
  \lean{QOU.MassDerivation.mu_e_ratio_approx}
  \uses{def:mu-e-ratio}
  $|m_\mu/m_e - 206.768| < 10^{-3}$.
\end{lemma}

% ── Quantum Planck Approximation ──────────────────────────────────
\begin{lemma}[Quantum Planck Approximation]
  \label{lem:hq-approx}
  \lean{QOU.MassDerivation.quantum_planck_approx}
  \uses{def:quantum-planck, def:substrate-q}
  $|\hbar_q(q_0) - 0.09908| < 10^{-4}$.
\end{lemma}

% ── Derive Substrate q ────────────────────────────────────────────
\begin{theorem}[Derivation of the Substrate Parameter]
  \label{thm:derive-substrate-q}
  \lean{QOU.MassDerivation.derive_substrate_q}
  \uses{conj:topological-mass, def:vol-figure-eight, def:mu-e-ratio,
        def:quantum-planck}
  Given the topological mass conjecture and experimental data,
  the substrate parameter $q$ must satisfy $|q - 1.1097| < 0.0001$.

  Derivation:
  \begin{enumerate}
    \item From $m_\mu/m_e = \mathrm{Vol}(4_1)/\hbar_q^2$, solve for
          $\hbar_q^2 \approx 2.02988/206.768 \approx 0.009817$.
    \item $\hbar_q \approx 0.09908$.
    \item From $\hbar_q = 1 - 1/q$, invert: $q = 1/(1-\hbar_q) \approx 1.1097$.
  \end{enumerate}
\end{theorem}

% ── Mass Ratio from Volume ────────────────────────────────────────
\begin{lemma}[Mass Ratio from Volume]
  \label{lem:mass-ratio-from-volume}
  \lean{QOU.MassDerivation.mass_ratio_from_volume}
  \uses{def:vol-figure-eight, def:mu-e-ratio, def:quantum-planck}
  $\hbar_q \neq 0 \implies \mathrm{Vol}(4_1)/\hbar_q^2 = m_\mu/m_e
  \implies \hbar_q^2 = \mathrm{Vol}(4_1)/(m_\mu/m_e)$.
\end{lemma}


\chapter{Calculations Module (Registry Bridge)}

% ── Registry Volume Consistency ───────────────────────────────────
\begin{theorem}[Registry Volume Consistency]
  \label{thm:registry-volume}
  \lean{QOU.Calculations.registry_volume_eq_mass_derivation}
  \leanok
  \uses{def:figure-eight, def:vol-figure-eight}
  The figure-eight knot's registry volume equals the
  \texttt{MassDerivation} volume constant.
\end{theorem}

% ── Skein Mass Residue ────────────────────────────────────────────
\begin{theorem}[Skein-Deformed Mass Residue]
  \label{thm:skein-mass-residue}
  \lean{QOU.Calculations.skein_mass_residue}
  \uses{thm:registry-volume, def:mu-e-ratio, def:quantum-planck}
  The muon-electron mass ratio equals
  $\mathrm{Vol}(4_1)/\hbar_q^2$ where the volume is sourced from
  the Alexander-Briggs knot registry.
\end{theorem}

% ── Substrate Error Bound ─────────────────────────────────────────
\begin{theorem}[Substrate Error Bound]
  \label{thm:substrate-error-bound}
  \lean{QOU.Calculations.substrate_error_bound}
  \uses{thm:derive-substrate-q, thm:registry-volume}
  $|q_{\mathrm{derived}} - 1.1097| < 0.0001$.
\end{theorem}

% ── Full Derivation Chain ─────────────────────────────────────────
\begin{theorem}[Full Derivation Chain]
  \label{thm:full-derivation-chain}
  \lean{QOU.Calculations.full_derivation_chain}
  \uses{def:figure-eight, def:mu-e-ratio, thm:substrate-error-bound}
  End-to-end witness:
  \[
    \mathrm{KnotRegistry}(4_1).\mathrm{vol}
      \;\longrightarrow\; \hbar_q^2
      \;\longrightarrow\; \hbar_q
      \;\longrightarrow\; q
      \;\longrightarrow\; |q - 1.1097| < 0.0001
  \]
\end{theorem}

% ── CODATA Consistency ────────────────────────────────────────────
\begin{lemma}[$\hbar_q^2 > 0$]
  \label{lem:hq-squared-pos}
  \lean{QOU.Calculations.h_q_squared_pos}
  \uses{def:vol-figure-eight, def:mu-e-ratio}
  Both $\mathrm{Vol}(4_1)$ and $m_\mu/m_e$ are positive,
  so $\hbar_q^2 = \mathrm{Vol}/\mathrm{ratio} > 0$.
\end{lemma}

\begin{lemma}[$\hbar_q \in (0,1)$]
  \label{lem:hq-unit-interval}
  \lean{QOU.Calculations.h_q_value_in_unit_interval}
  \uses{lem:hq-squared-pos}
  Ensures $q > 1$, as required by the quantum group framework.
\end{lemma}

\begin{lemma}[$q > 1$]
  \label{lem:derived-q-gt-one}
  \lean{QOU.Calculations.derived_q_gt_one}
  \uses{lem:hq-unit-interval}
  The derived substrate parameter satisfies $q > 1$
  ($q = 1$ would collapse to the classical limit).
\end{lemma}
